% LaTeX テンプレート — 授業資料用 (LuaLaTeX を想定、両面: 奇数ページにヘッダー)
\documentclass[11pt,twoside]{article}

% LuaLaTeX 向け設定: 日本語対応に luatexja-fontspec と fontspec を使用
\usepackage{fontspec}
\usepackage{luatexja-fontspec}

% 必要に応じて日本語フォントを明示（環境に応じて変更してください）
% 例: \setmainjfont{Noto Serif CJK JP}

% 単位表示の統一: siunitx を使用
\usepackage{siunitx}
\sisetup{detect-mode=true, per-mode=symbol}

% その他のパッケージ
\usepackage{tikz}
\usepackage{fancybox}
\usepackage{ascmac}
\usepackage{amsmath}
\usepackage{geometry}
\usepackage{fancyhdr}
\usepackage{lipsum} % サンプル用ダミーテキスト
\usepackage{ulem}
\usepackage{enumitem}
\usepackage{amssymb}
\usepackage{dashrule}

%Tikz
\usepackage{physmotion}

% 使い方: 以下を編集してください
\newcommand{\SubjectTitle}{水平投射・斜方投射} %テーマ
\newcommand{\ClassRoomInfo}{\underline{\hspace{0.5cm}}年\underline{\hspace{0.5cm}}組　氏名\underline{\hspace{3cm}}}

% 余白設定
\geometry{top=20mm,bottom=25mm,left=20mm,right=20mm,heightrounded}

% ヘッダー設定: 奇数ページに表示（左: 科目と内容、右: 学年組欄）
\setlength{\headheight}{14pt}
\setlength{\headsep}{36pt}
\pagestyle{fancy}
\fancyhf{} % 既定のヘッダー/フッターをクリア
\fancyhead[LO]{\small 物理基礎「\SubjectTitle」} %Odd page Left headerを設定
\fancyhead[RO]{\small \ClassRoomInfo} % Odd page Right headerを設定
\renewcommand{\headrulewidth}{0.4pt}
\renewcommand{\footrulewidth}{0pt}
% フッター中央にページ番号を表示
\fancyfoot[C]{\small\thepage}

% 偶数ページはヘッダーを空にして、本文が上から始まるようにする
\fancyhead[LE,CE,RE]{}

% ドキュメント開始
\begin{document}
%=====================================================================================
\section{\textbf{復習}}
\begin{itembox}[l]{\textbf{〇等加速度直線運動}}
	\begin{minipage}[t]{0.70\linewidth}
	鉛直投げ下ろし運動の変位・速度の式
	\Large
	\[
	\left(
	\begin{array}{ll}
		\text{①} \quad v = & \\
		\text{②} \quad x = & \\
		\text{③} \quad v^2 - v_0^2 = 2ax &
	\end{array}
	\right.
	\]
	\begin{center}
		% \motiondiagram{初速度}{終速度}{加速度}：画像を挿入
		\motiondiagram{v_0}{v}{a}{t}
	\end{center}	
	\end{minipage}
	\hfill
	\begin{minipage}[t]{0.26\linewidth}
		\begin{shadebox}
			\begin{align*}
				v\,[\si{\metre\per\second}] &\text{：速度} \\
				v_0\,[\si{\metre\per\second}] &\text{：初速度} \\
				a\,[\si{\metre\per\second\squared}] &\text{：加速度} \\
				t\,[\si{\second}] &\text{：時間} \\
				x\,[\si{\metre}] &\text{：変位} \\
			\end{align*}
		\end{shadebox}
	\end{minipage}
\end{itembox}\\
%=====================================================================================
\begin{itembox}[l]{\textbf{〇鉛直投げ下ろし・鉛直投げ上げ}}
	\begin{minipage}[t]{0.70\linewidth}
		鉛直投げ下ろし運動の変位・速度の式
		\Large
		\[
		\left(
		\begin{array}{ll}
			\text{①} \quad v = & \\
			\text{②} \quad y = & \\
			\text{③} \quad v^2 - v_0^2 = 2gy &
		\end{array}
		\right.
		\]
		\begin{shadebox}
			$g\,[\si{\metre\per\second\squared}] \text{：重力加速度}$
			$v_0=0$で自由落下
		\end{shadebox}
	\end{minipage}
	\hfill
	\begin{minipage}[t]{0.26\linewidth}
		\begin{center}
			% \throwdowndiagram{初速度v_0}{終速度v}{加速度g}{時刻変数t}：画像を挿入
			\throwdowndiagram{v_0}{v}{g}{t}
		\end{center}
	\end{minipage}
\end{itembox}


\clearpage
%=====================================================================================
% 水平投射
\section{水平投射}
	% \lipsum[1] //ダミーテキスト
	\begin{minipage}[t]{0.34\linewidth}
		\projectilediagram{v_0}{v_x}{v_y}{v}{g}{t}
		\\\\\\
		\begin{itembox}[l]{補足}
			I. 実際の速度$v$[m/s]は\\\qquad
			\(
				v = 
			\)\\
			で求めることができる。\\
			\to 詳しくはベクトルを勉強\\\\
			II. ボールの軌道の導出\\
			\to 方針：右の(i),(ii)式から$t$を消して、$x$と$y$と定数$g$,$v_0$のみの式にする。\\\\\\\\\\\\\\
		\end{itembox}
	\end{minipage}
	\hfill
	\begin{minipage}[t]{0.58\linewidth}
		
		\ovalbox{①水平投射とは？}\\
		左図のような重力($y$軸方向)がかかる運動のうち、初速度$v_0$が水平($x$軸方向)成分のみである物体の運動。\\\\
		\ovalbox{②今までとの違い}\\
		重力方向以外にも運動する。\\
		　\to \textbf{x軸方向も考えなければならない}\\\\
		\ovalbox{③セットアップ}
			\begin{itemize}[leftmargin=2em,itemsep=0pt,topsep=2pt]
				\item[(i)]
				軸の向き：
				初速度の向きを $x$ 軸の正、
				重力加速度の向きを $y$ 軸の正とする。

				\item[(ii)]
				$t=0$ のとき、
				物体は原点 $(x,y)=(0,0)$ にある。

				\item[(iii)]
				重力加速度の大きさを $g$、
				初速度の大きさを $v_0$ とする。
			\end{itemize}
			　\\
		\ovalbox{④$t$秒後の物体の位置と速度}　\uuline{\textbf{大事！！}}\\
		\textbf{point!：}$x$と$y$を分けて考える\\
		\textbf{point!：}$y$軸方向のみ加速している
		\begin{itemize}[leftmargin=2em,itemsep=0pt,topsep=2pt]
			\item[(i)]
			$x$軸方向の運動：\\
			\(
			\left(
			\begin{array}{ll}
				\text{①} \quad a_x = & \\
				\text{②} \quad v_x = & \\
				\text{③} \quad x = & \\
			\end{array}
			\right.
			\)\\
			\to x軸方向の運動は\underline{\hspace{3.5cm}}と一致！\\
			\item[(ii)]
			$y$軸方向の運動：\\
			\(
			\left(
			\begin{array}{ll}
				\text{①} \quad a_y = & \\
				\text{②} \quad v_y = & \\
				\text{③} \quad y = & \\
			\end{array}
			\right.
			\)\\
			\to y軸方向の運動は\underline{\hspace{3.5cm}}と一致！
		\end{itemize}
	\end{minipage}
\newpage
%=====================================================================================
	\section{演習問題(水平投射)}
	\begin{shadebox}
		問題を解くときのポイント：\\
			①軸の向きを意識する(軸と同じ方向のベクトルは+、逆向きのベクトルは-になる)\\
			②x軸方向とy軸方向を分けて考える\\
			③公式を使って、ほしい量を求める\\
			　ただし、この時どの文字が最後まで残っていいか意識する\\
			　例)前ページの軌道を計算する際には、$y$を$x$の式で表すため、$t$は削除する。\\
			　　\,\,\,\,\,$g$や$v_0$は問題文で与えられる初期条件なので残してもよい\\
			　　\,\,\,\,\,問題によっては「$y$を$x$，$g$，$v_0$を用いて表せ」などと指示されることも多い\\
			③単位を確認して、軽い確認ができる\\
			　例)速度を求める問題で、答えの単位が\si{\metre\per\second$^2$}になっていたら間違っている可能性が高い
	\end{shadebox}
	　\\
	問1：以下の図のような...
	\clearpage

	% 斜方投射
	\section{斜方投射}
		\begin{minipage}[t]{0.44\linewidth}
			\obliqueprojectilediagram{v_0}{\theta}{v_{x0}}{v_{y0}}{g}{t}\\
			\begin{itembox}[l]{補足}
				I.初速度について(三角関数の復習)\\\\\\\\
				II. 最高点について\\\\\\\\
				III. ボールの軌道の導出\\
				\to 方針：右の(i),(ii)式から$t$を消して、$x$と$y$と$g$，$v_0$，$\theta$のみの式にする。\\\\\\\\\\\\\\\\\\
			\end{itembox}
		\end{minipage}
		\hfill	
		\begin{minipage}[t]{0.52\linewidth}
			\ovalbox{①斜方投射とは？}\\
			左図のような重力がかかる運動のうち、初速度$v_0$が任意の方向を向いている物体の運動。\\\\
			\ovalbox{②今までとの違い}\\
			$v_0$が$x$成分と$y$成分の両方の値を持ちうる。\\\\
			\ovalbox{③セットアップ}
				\begin{itemize}[leftmargin=2em,itemsep=0pt,topsep=2pt]
					\item[(i)]
					軸の向き：
					水平右向きを $x$ 軸の正、
					鉛直上向きを $y$ 軸の正とする。
	
					\item[(ii)]
					$t=0$ のとき、
					物体は原点 $(x,y)=(0,0)$ にある。
	
					\item[(iii)]
					重力加速度の大きさを $g$、
					初速度の大きさを $v_0$、仰角を $\theta$ とする。
	
					\item[(iv)]
					初速度の補足：各方向の初速度は以下のようになる。\\
					$v_{x0} = v_0\cos\theta$（$x$軸方向の初速度）\\
					$v_{y0} = v_0\sin\theta$（$y$軸方向の初速度）\\
				\end{itemize}
			\ovalbox{④$t$秒後の物体の位置と速度}　\uuline{\textbf{大事！！}}\\
			\textbf{point!：}$x$と$y$を分けて考える\\
			\textbf{point!：}$y$軸方向のみ加速している
			\begin{itemize}[leftmargin=2em,itemsep=0pt,topsep=2pt]
				\item[(i)]
				$x$軸方向の運動：\\
				\(
				\left(
				\begin{array}{ll}
					\text{①} \quad a_x = & \\
					\text{②} \quad v_x = & \\
					\text{③} \quad x = & \\
				\end{array}
				\right.
				\)\\
				\to x軸方向の運動は\underline{\hspace{3.5cm}}と一致！\\
				\item[(ii)]
				$y$軸方向の運動：\\
				\(
				\left(
				\begin{array}{ll}
					\text{①} \quad a_y = & \\
					\text{②} \quad v_y = & \\
					\text{③} \quad y = & \\
				\end{array}
				\right.
				\)\\
				\to y軸方向の運動は\underline{\hspace{3.5cm}}と一致！
			\end{itemize}
	\end{minipage}\\\\\\
	%=====================================================================================
	\begin{shadebox}
		まとめ\\\\\\\\\\\\\\\\\\\\
	\end{shadebox}
	\clearpage

	\section{演習問題(斜方投射)}
	\clearpage
\end{document}

